\documentclass[11pt,a4paper]{article}

% ------------------------------------------------------------------ setup
\usepackage[T1]{fontenc}
\usepackage[utf8]{inputenc}
\usepackage{newpxtext,newpxmath}
\usepackage[scaled=0.9]{inconsolata}
\usepackage[margin=2.4cm]{geometry}
\usepackage[english]{babel}
\usepackage{microtype}
% No "--" ligature in code (run --js).
\DisableLigatures[-]{family = tt*}
\usepackage{booktabs}
\usepackage{longtable}
\usepackage{array}
\usepackage{tabularx}
\usepackage{enumitem}
\usepackage[table]{xcolor}
\usepackage{xurl}
\usepackage[style=numeric-comp,sorting=nyt,maxbibnames=8,minbibnames=6,backend=biber,isbn=false,giveninits=false]{biblatex}
\usepackage{csquotes}
\usepackage[normalem]{ulem}
\usepackage{hyperref}
\usepackage{titlesec}
\usepackage{fancyhdr}

\addbibresource{literature-review.bib}
% The addendum names the register card each source rests on; print it small.

\DeclareFieldFormat{addendum}{\footnotesize #1}

\definecolor{swblue}{HTML}{1F4E79}
\definecolor{swgreen}{HTML}{2E6B30}
\definecolor{swred}{HTML}{8B1A1A}
\definecolor{swgrey}{HTML}{555555}
\hypersetup{
  colorlinks=true,
  linkcolor=swblue,
  citecolor=swgreen,
  urlcolor=swred,
  pdftitle={Ideas to Steal: A Literature Review for Stitchwort},
  pdfauthor={Stitchwort project, SW-894},
  pdfsubject={Summary of the SW-894 research registers},
}

\titleformat{\section}{\Large\bfseries\color{swblue}}{\thesection}{0.8em}{}
\titleformat{\subsection}{\large\bfseries\color{swblue}}{\thesubsection}{0.7em}{}
\titleformat{\subsubsection}{\normalsize\bfseries}{\thesubsubsection}{0.6em}{}

\pagestyle{fancy}
\fancyhf{}
\fancyhead[L]{\small\color{swgrey}Stitchwort literature review (SW-894)}
\fancyhead[R]{\small\color{swgrey}30 October 2026}
\fancyfoot[C]{\small\thepage}
\renewcommand{\headrulewidth}{0.4pt}

\setlist{itemsep=2pt,topsep=3pt,parsep=0pt}
\setlength{\parskip}{3pt}
\setlength{\emergencystretch}{2em}

% Code and paths. \path breaks at punctuation and needs no escaping.
\newcommand{\code}[1]{\texttt{#1}}
\newcommand{\file}[1]{\path{#1}}
\newcommand{\tkt}[1]{SW-#1}
\newcommand{\dec}[1]{D#1}
% Claim labels, used sparingly, as in the registers.
\newcommand{\inferred}{\textsc{[inferred]}}
\newcommand{\unverified}{\textsc{[unverified]}}
\newcommand{\run}{\textsc{[run]}}

\newcolumntype{L}[1]{>{\raggedright\arraybackslash}p{#1}}
\renewcommand{\arraystretch}{1.15}

\title{\color{swblue}\textbf{Ideas to Steal}\\[4pt]
  \large A literature review for Stitchwort, summarising the seven SW-894 research registers}
\author{Stitchwort project, board ticket SW-894}
\date{30 October 2026}

\begin{document}
\maketitle
\thispagestyle{empty}

\begin{abstract}
\noindent
Stitchwort is a small, total, dependently typed functional language written in Rust, with
algebraic effects and a small trusted kernel. Between 29 and 30 October 2026 it was the subject
of a literature review in seven clusters: effects and handlers, elaboration, proof automation,
the kernel and type theory, inductive types and recursion, infrastructure and trust, and
compilation and runtime. The window was 2022 onwards. Across the seven registers 449 sources were
screened and 160 carded; 153 of the cards rest on a full text read for the review and 7 were not
read. A first attempt on 29 October wrote several registers from abstracts. Every register was
then redone or audited against full texts on 30 October, and this summary cites only sources
read in full.

The review has already changed Stitchwort. Porting a 2026 kernel-bug shape from Lean4Lean
found a real soundness bug, an unchecked universe-level overflow, which was fixed (SW-902).
Fourteen further must-reject programs, shaped like 2026 bugs in other provers, are all rejected.
The confluence argument for the extra \code{Nat.add} rules now rests on the Taming of the Rew.
Two kernel hazards are recorded as decisions against future work (SW-93, SW-55). The rest of the
findings point the same way. The literature broadly confirms Stitchwort's architecture: an
untrusted elaborator, certificates checked by the kernel, features fixed when a declaration is
admitted. What it adds is mostly sharper tests and oracles, a handful of kernel restrictions,
proof automation that stays inside the elaborator, and one elaborator redesign (a single
postponement queue). Section~\ref{sec:recs} gives a prioritised list and
Section~\ref{sec:decisions} the decisions still open.
\end{abstract}

\tableofcontents
\clearpage

% =====================================================================
\section{Stitchwort in brief}\label{sec:brief}

This section is for a reader without access to the repository. Everything here is from the
project's own documents (\file{README.md}, \file{DESIGN.md}, \file{docs/core.md}) as they stand on
\code{main} on 30 October 2026, and from the registers' descriptions of the code.

\paragraph{The language.} Stitchwort is a strict, pure, dependently typed functional language,
implemented in Rust (the \file{crates/} directory: \code{syntax}, \code{elab}, \code{core},
\code{rt}, \code{js}, \code{native}, \code{driver}, \code{cli} and others). It combines four
ideas that usually live in separate languages: functions terminate by default, and anything
that might not says so in its type; dependent types let programs carry proofs about themselves;
effects are tracked in types and implemented by handlers, in direct style; and definitions are
identified by a hash of their contents, so tests and proofs whose dependencies haven't changed
are never re-run. \code{Prop} is impredicative and erased at run time; universes are numbered
and not cumulative; the kernel has axiom K. The one axiom is \code{Quot.sound}, and
\code{stitchwort axioms} reports whether a definition depends on it.

\paragraph{Abilities.} Effects are called \emph{abilities}, after Unison. A computation with
abilities is translated into data, \code{Effect.Comp}, with constructors \code{Done} and
\code{Call}; a handler is a fold over it (\code{Effect.foldComp}), and sequencing is
\code{Effect.bind} (\file{crates/elab/src/std/core/Effect.sw:12-28}). Handlers are deep and
resumptions may be used any number of times (decision D1162). Because handlers are ordinary
folds, the kernel's own positivity and termination checks prove that handled programs
terminate; what remains trusted is that the translation is faithful to the source semantics.

\paragraph{The kernel.} A small trusted kernel (\file{crates/core}) re-checks every declaration
the elaborator produces. It has case trees with size-change termination rather than
recursors, an untyped conversion checker with eta for functions, two extra conversion rules for
\code{Nat.add} (R3/R4, pinned by content hash), and a short list of natives (fast
implementations of standard definitions). The elaborator's unifier is untyped and may be
incomplete; that is safe because the kernel re-checks the result.

\paragraph{The elaborator.} The elaborator (\file{crates/elab}) does pattern unification with
pruning, keeps glued calls (definitions compared by arguments before unfolding), infers
implicit arguments, solves instance obligations in a per-declaration loop without
backtracking, and has records with literals, projections and a field-name rule for literals
whose type is unknown. It postpones lambda and method-section arguments within one application,
and since SW-896 puts aside unification pairs stuck on an unknown.

\paragraph{Proof tools.} Proofs are terms, written with \code{match}, \code{J}, \code{rewrite},
and a few automation words: \code{rfl}; \code{cases} (case splits with closers: computation,
constructor clash, injectivity, a hypothesis, \code{omega}); \code{omega} and \code{ring}, which
search outside the kernel and produce a certificate the kernel checks; \code{auto}, which uses
\code{Decidable} instances and hypotheses in scope; and \code{decreasing} for recursion on a
measure, with \code{Nat} fuel rather than \code{Acc}. A counterexample search runs on a failing
theorem, filters random inputs by its hypotheses, and is advisory only. A \code{?} hole leaves a
goal and checking continues.

\paragraph{Compilation.} Programs run on a fast runtime (\file{crates/rt}): an erased IR,
register bytecode with moves on last use and Perceus/Lean-style reuse tokens, small numbers
inline with a bignum fallback, tail recursion modulo context by destination passing, and state
handlers run as loops.
Native code is generated with Cranelift (\file{crates/native}) and keeps the machine's reference
counting decisions. A JavaScript backend serves \code{run --js} and the extraction of proven code
as an ES module with TypeScript types. A reference evaluator (\code{run --reference}) runs
programs on the type checker's own evaluator. A browser playground runs the machine as
WebAssembly.

\paragraph{The corpus.} A standard library (\code{std}); examples, including must-reject
examples in \file{examples/soundness/} (107 of them after 30 October); Advent of Code 2024 and
2025 with proofs (\file{aoc/}); and a roguelike, \file{games/rogue}, with its own proof notes.
Tests run in tiers (lane, merge, daily, full), with fuzzers for the parser, formatter,
elaborator and kernel, and differential tests across the four ways of running a program.

% =====================================================================
\section{Method}\label{sec:method}

\subsection{The seven clusters}

The plan, in SW-894, divided the field into seven clusters, each with board hooks:
\begin{enumerate}[label=\arabic*.]
  \item \textbf{Kernel and type theory}: rewrite rules, definitional proof irrelevance, record eta,
        observational type theory, quotients, universe polymorphism.
  \item \textbf{Inductive types and recursion}: nested and inductive-inductive types, size-change
        termination, well-founded recursion.
  \item \textbf{Elaboration}: unification, postponement, implicit arguments, glued evaluation,
        error messages.
  \item \textbf{Proof automation}: rewriting, decision procedures, e-graphs, counterexample search.
  \item \textbf{Effects and handlers}: rows, evidence and capability passing, handler
        termination, effects with dependent types.
  \item \textbf{Compilation and runtime}: erasure, reference counting and reuse, interpreters,
        JS and WebAssembly.
  \item \textbf{Infrastructure and trust}: content-addressed code, incremental compilation,
        testing and fuzzing of checkers, agent-driven development.
\end{enumerate}
Each cluster has one register in \file{docs/research/}, with the same parts: method and queries,
screened list, cards (about five lines per kept source), deep reads of the top three to five,
ranked ``steal next'', proposed tickets, and at most a few decisions for Matt.

\subsection{Window}

Matt set the window at 2022 onwards (D1831), with snowballing inside it. Older work appears only
where in-window work or Stitchwort's own code relies on it, and is marked \emph{classic} in the
registers and in this summary (for example the Taming of the Rew, 2021, and Counting Immutable
Beans, 2019).

\subsection{The reading protocol}

After the failure described next, every register follows one protocol:
\begin{itemize}
  \item \textbf{Full text or nothing.} PDFs were read with \code{pdftotext} from
        \code{nix shell nixpkgs\#poppler-utils}; source files, manual pages and issues were read
        whole. For a long theory paper, ``read'' means the prose, definitions, examples,
        evaluation and conclusions, with proofs skimmed; each card names the sections read.
  \item \textbf{Counts first.} Every register opens with screened, carded, read in full and not
        read.
  \item \textbf{NOT READ supports nothing.} Every card says ``Full text: read (source)'' or
        ``NOT READ''. A card not read cannot be deep-read, ranked or used for a recommendation.
  \item \textbf{The paywall rule.} Sources were tried in order (arXiv, authors' pages, DBLP,
        Semantic Scholar, Dagstuhl, HAL, ACM, the Internet Archive). A full text that couldn't be
        reached goes in a ``for Matt to fetch'' table only if it would change something; the
        rest are listed as not reached.
  \item \textbf{No personal data} was sent to any service.
\end{itemize}

\subsection{The 29 October failure and the 30 October redo}

The first attempt, on 29 October, went badly; SW-897 records how. No PDF tool was installed and
the web-fetch tool could not read PDFs, and nobody checked before the lanes started, so the
elaboration, proof-automation and kernel registers were written largely from abstracts. Each
lane listed its unread sources at the end of its report, and the orchestrator relayed the
findings and dropped the caveat. Up to four research lanes ran at once and exhausted the
session's shared web-search pool, and one lane's claim that the budget had run out was relayed
without checking. The orchestrator once sent Matt's email address to an open-access API. The
kernel redo's readers later found that the first kernel register misstated four papers.

Matt paused the work (SW-897) and restarted it on 30 October one lane at a time (D1974), with the
protocol above, a tested \code{pdftotext} toolchain, reserved decision numbers and a planned
search budget. The elaboration, proof-automation and kernel registers were redone from full
texts; the effects, inductives, infrastructure and compilation registers were audited card by
card against full texts. Each redo also lists the 29 October claims it found wrong: five in
effects, six each in elaboration, proof automation and kernel. The audits found one wrong card
each in effects (a steal), infrastructure (numbers taken from an anonymous submission copy) and
compilation (a thesis example attributed to Lean's library), and none in inductives, though one
of its assumptions about Stitchwort was wrong.

\subsection{Counts}

Table~\ref{tab:counts} gives each register's counts as its first line states them. Some sources
are carded in two registers (for example Lean's functional induction, the Hazel error-recovery
paper and the PLDI 2022 inductive-relations paper), so the totals count cards, not distinct
works.

\begin{table}[ht]
\centering
\small
\caption{Counts per cluster, from each register's first line (30 October 2026).}
\label{tab:counts}
\begin{tabular}{@{}L{3.2cm}rrrrL{5.6cm}@{}}
\toprule
Cluster & Screened & Carded & Read & Not read & How checked \\
\midrule
Effects and handlers & 75 & 11 & 11 & 0 & Audited: 6 OK, 4 minor, 1 wrong (fixed). The 11 read are 10 papers and Lean's manual (C3). \\
Elaboration & 64 & 20 & 20 & 0 & Redone from full texts; 6 old claims wrong. \\
Proof automation & 72 & 29 & 25 & 4 & Redone from full texts; 6 old claims wrong. \\
Kernel and type theory & 62 & 27 & 26 & 1 & Redone from full texts; 6 old claims wrong. \\
Inductive types and recursion & 56 & 21 & 21 & 0 & Audited (19 cards): 11 OK, 8 minor, 0 wrong; 2 cards first NOT READ were read. \\
Infrastructure and trust & 68 & 29 & 28 & 1 & Audited (24 cards): 18 OK, 5 minor, 1 wrong; searching finished. \\
Compilation and runtime & 52 & 23 & 22 & 1 & Audited (21 cards): 12 OK, 8 minor, 1 wrong; C12 read at the audit. \\
\midrule
\textbf{Total} & \textbf{449} & \textbf{160} & \textbf{153} & \textbf{7} & \\
\bottomrule
\end{tabular}
\end{table}

Web searches were few in the redo: five each for elaboration and kernel, six for proof
automation, one for the effects audit, none for the inductives audit and twelve for the
infrastructure audit, which finished the searching its 29 October version had been cut short of.

% =====================================================================
\section{What has already been acted on}\label{sec:acted}

Table~\ref{tab:acted} lists what the review has already changed, with the commits and tickets
that did it. Each item was checked in \code{git log} and on the board on 30 October.

\begin{longtable}{@{}L{4.4cm}L{7.0cm}L{3.8cm}@{}}
\caption{Already acted on.}\label{tab:acted}\\
\toprule
Item & What was done & References \\
\midrule
\endfirsthead
\toprule
Item & What was done & References \\
\midrule
\endhead
\bottomrule
\endfoot
Universe-level overflow (soundness bug) &
Found by the lane porting 2026 kernel-bug shapes (kernel register T1), on Lean4Lean's
\code{looseBVarRange} overflow shape \cite{carneiro2026lean4lean}. \code{def wrap: Type 0 := Type
4294967295} was accepted: levels were \code{u32} and the step to the next universe wrapped. Fixed:
\code{MAX\_LEVEL} is 65535 and \code{Sort::succ} returns no successor for it
(\file{crates/core/src/term.rs:33}, \file{:40}); parser, elaborator and cache limits; five
\file{universe-*.sw} examples and \file{crates/core/tests/universe_levels.rs}. &
SW-902 (D1984); e1d444d7, da80af16, d8d36404, d3c2dee4; merged 65f6a236. \\
\addlinespace
14 must-reject tests for 2026 kernel-bug shapes &
Programs for Rocq CVE-2026-72703 and 72705, Arena \code{k-rec-conv},
\code{eta-unit-like-proof-field} and \code{forged-quot-ind}, and Lean's order-dependent cache
(lean4\#14806); all rejected (\code{E-TERMINATION}, \code{E-REFL}, \code{E-DUPLICATE}), about 0.13\,s
of CPU each. A kernel-level test builds the same terms by hand
(\file{crates/core/tests/kernel_bugs_2026.rs}, 6/6). &
Kernel register T1 (D1983); 3208ab72, 1ae69c78; merged f38544f6. \\
\addlinespace
\code{+} confluence on the Taming of the Rew &
\file{docs/proposals/nat-add-recursion.md} \S9.4 now argues confluence of R1--R4 by the triangle
property \cite{cockx2021taming} instead of Newman's lemma; a new assumption A1$'$ (the core's own
triangle property); test \code{the\_parallel\_rule\_is\_reached}
(\file{crates/driver/tests/nat_add_rule.rs:160}) checks that the evaluator reaches the parallel
rule's right side. &
Kernel register T4 (D1983); c1084edd; merged f9521afe. \\
\addlinespace
Recursive propositions (SW-93 warning) &
D1980 (Matt): refuse eliminating a self-referential proposition into data, as part of SW-93,
after measuring that nothing in the corpus does it. The warning is in SW-93's body. &
D1980; ba9ba8a7. \\
\addlinespace
Record eta reminders &
D1981 (Matt): once proofs are irrelevant, record eta only for records with a field that isn't a
proof. Reminders placed on SW-93, SW-55, SW-789 and SW-568. &
D1981; ba9ba8a7. \\
\addlinespace
Counterexample search misses &
SW-903 records that \code{check --counterexamples} misses a false theorem silently when its
hypothesis holds only for small inputs; re-run by the orchestrator on \code{main}. &
SW-903 (backlog); 9094029f. \\
\addlinespace
Congruence closure in \code{cases} &
SW-904 records the idea. D2000 (Matt): not prioritised yet; first survey existing proofs for
steps it would remove. &
SW-904 (backlog), D2000; 2a66a5f9. \\
\addlinespace
One postponement queue &
D1990: a design note, \file{docs/proposals/postponement.md}, comparing Lean 4, Agda, Idris~2, Rocq
and OmniML. D2012 (Matt): design for the best end state, so draft 2 makes one queue in the
single elaboration pass, with prototype measurements and a brief for an external reviewer. The
review has been received; draft 3 (D2015) is revising the note. &
SW-817 (doing); 7b751873 / 8cd3cdd6; c35e458b / bb9e432a; 3d0036f8. \\
\addlinespace
SW-805, field-name rule &
D1991: keep the rule as a name-resolution rule; it was merged as a last resort, used only where
nothing else gives the literal's type. &
SW-805 (done); 8f52025b; merged afea8395. \\
\addlinespace
\file{decide.md} (D1971--D1973) &
Design note for SW-800 and SW-807: extend \code{cases} with a decision closer and a split of a
finite variable, derive a \code{Finite} interface, never \code{native\_decide}. Merged; the
proof-automation register finds that D1971 stands, and the inductives audit withdrew its D1912
in favour of D1972. The decisions await a review artefact. &
88e9104b, 0ddc8e2c; SW-800 (next). \\
\addlinespace
Handler termination in \file{DESIGN.md} &
Effects register T1 (D1852): DESIGN.md described handler termination as a tracked conjecture;
it now says the kernel checks it and names the two trusted items. &
SW-899; 4e33920c. \\
\end{longtable}

% =====================================================================
\section{Effects and handlers}\label{sec:effects}

\textbf{Register:} \file{docs/research/effects.md}. 75 screened, 11 carded, all read; audited.

\subsection{Key findings}
\begin{itemize}
  \item \textbf{Monad laws by conversion have a price.} Yang and Wu get the monad laws
    judgementally for handlers carried by lawless raw monads, by giving up the law that
    handling commutes with \code{let}, since with lawless monads it ``would result in
    inconsistency'' (Remark 2.11) \cite{yang2026judgemental}. Their calculus is System F$\omega$
    with effects, not dependently typed, and its \S6 leaves open whether a dependently typed
    language can have all the laws judgementally. Their model is a CPS encoding in which the
    laws hold by $\beta$.
  \item \textbf{Hoare triples that yield equations.} Kidney, Yang and Wu's H-elim turns a
    triple into a program equation that can be substituted in any context
    \cite{kidney2024hoare}. Its hard part comes from quotients; on Stitchwort's unquotiented
    \code{Comp} it is an easy induction, but it still needs \code{Quot.funext} at \code{Call}
    and, the audit found, it helps neither rogue's ``gap 2'' nor SW-548.
  \item \textbf{Specifications as predicate transformers.} Lean's \code{Std.Do},
    \code{mvcgen} and \code{vcgen} give each program a weakest-precondition reading and reduce
    triples to verification conditions with no weakest preconditions left
    \cite{lean2026mvcgen}. \code{vcgen} accepts any program type with a \code{WP} instance,
    including syntax trees.
  \item \textbf{Declared fast paths.} Lexa declares handlers tail-resumptive or abortive and
    runs them in place, ``reducing running-time complexity from quadratic to linear in some
    cases'' \cite{ma2024lexa}; it treats multi-shot resumptions as a non-goal. A per-effect
    strategy chosen by invocation rate gives exceptional effects zero overhead
    \cite{ma2025zero}.
  \item \textbf{TRMC through effectful callbacks.} Koka keeps an effectful \code{map} in
    constant stack either by copying a shared context or by running a second version whenever a
    multi-shot handler is in scope \cite{leijen2023trmc}.
  \item \textbf{Termination.} Hamana proves a handler strongly normalising by the argument
    the kernel uses, that \code{M[x]} is smaller than \code{get(x.M[x])} \cite{hamana2022modular}.
  \item Smaller points: messages should never show a hidden effect variable
    \cite{tang2025modal}; Koka closes rows at \code{let}, a related but different rule from
    Stitchwort's D1124 \cite{ikemori2022closed}; interaction-tree specifications are a
    further-off fit, for coinductive programs \cite{silver2023itree}; Lean's \code{do}
    translation targets transformers, later replaced by CPS versions for speed
    \cite{ullrich2022do}.
\end{itemize}

\subsection{What Stitchwort could take}
The register's ranking, by value against effort (it was not re-ranked at the audit):
\begin{enumerate}
  \item \sout{Retire the handler-termination conjecture text} \emph{(done, SW-899)}.
  \item \code{Effect.bindCongAlways}, now narrowed to a small std lemma for rewriting a
        continuation under a fact the computation guarantees.
  \item A handler fast-path assertion or report (T5): \file{crates/rt/src/fusion.rs}'s
        recognisers return \code{Option} with no reason (\file{fusion.rs:242}, \file{:269},
        \file{:482}), so SW-861's slowdown (9.4\,s against 2.2\,s) passed silently. Re-measure
        SW-861 first: a later commit may already fuse its natural form \inferred.
  \item Handler specifications as predicate transformers (T4, D1851).
  \item The monad laws of \code{Effect.bind} and \code{Effect.foldComp} as trusted conversion
        rules (T3, D1850): a kernel change, design note only. Unlike R3/R4 these rules are
        higher-order and need funext, so the \code{Nat.add} confluence argument does not carry
        over; the note must weigh them against a CPS representation, which would lose induction
        on computations.
  \item Messages never show a bare row unknown.
  \item Guarded TRMC for an effectful \code{map} (SW-868 option (b$'$)).
  \item A per-ability runtime strategy (research).
  \item Modal effects instead of \code{\{e\}} on higher-order functions (research; conflicts
        with ``a plain arrow means pure and total'').
\end{enumerate}

\subsection{Acted on, and open}
SW-899 rewrote DESIGN.md's handler-termination text (D1852, closed). Open: D1850 and D1851
(Section~\ref{sec:decisions}).

% =====================================================================
\section{Elaboration}\label{sec:elab}

\textbf{Register:} \file{docs/research/elaboration.md}. 64 screened, 20 carded, all read; redone.

\subsection{Key findings}
\begin{itemize}
  \item \textbf{One postponement queue.} Lean postpones any elaborator that needs an unknown
    expected type (record instances, dot notation on an unknown type) and resumes them in a
    fixed ladder: resume while progress is made, then without postponing, then default
    instances, then report; it also has a body-first \code{let} \cite{lean2026elabsrc}. OmniML
    specifies the idea: a construct waits on the shape of one type, each unknown keeps a wait
    list, and a wait never discharged is an error \cite{obrien2025omniml}. ExEl runs every
    solver over one constraint store after each declaration \cite{liesnikov2023exel}.
    Stitchwort postponed only lambda and method-section arguments within one application
    (\file{crates/elab/src/check.rs:1903}, \code{apply\_args\_postponing}), so record literals
    (SW-817, SW-885) and let-bound lambdas (SW-847) failed at once.
  \item \textbf{Stitchwort's unifier is lossy for stuck calls.} Two forced calls of one
    definition with the same head have their spines unified as ordinary pairs, which may solve
    unknowns (\file{crates/elab/src/unify.rs:929-944}). That is the approximate rule: Coq and
    Lean use it by default, Agda only on request because it ``will cause Agda to ignore some
    possible solutions'' \cite{agda2026lossy}, and smalltt forbids it
    \cite{kovacs2021smalltt}. Polarity states the precise rule: congruence on arguments is safe
    for a neutral head, not a blocked one \cite{liesnikov2025polarity}. The kernel's re-check
    keeps it sound; nothing in \file{docs/core.md} says the elaborator may pick a different
    solution than the user meant.
  \item \textbf{SW-805's field-name rule is the one GHC kept.} GHC removed every
    type-directed field rule but kept ``only one datatype has all the fields'', because it is
    decided during name resolution \cite{gundry2020ghc366}; OmniML types a literal by the same
    closed-world reasoning \cite{obrien2025omniml}.
  \item \textbf{Errors.} In the one controlled study, correct hints raised success and speed
    for Agda newcomers and incorrect ones did no measurable harm (56 participants, 412
    submissions) \cite{khakimova2026hints}; flow-based messages showed no quantitative gain,
    and bare \code{?a} unknowns confused readers (117 participants kept)
    \cite{bhanuka2024flow}; total recovery needs a type for every erroneous term
    \cite{zhao2024marking}.
  \item \textbf{Benchmarks.} smalltt's \code{idTest} (40 nested \code{id}s) and \code{pairTest}
    (30 nested \code{let} pairs) take 0.000\,s in smalltt and time out in Agda, Coq, Lean and
    Idris; \code{vecTest} (960 elements) takes 0.078\,s against 0.244\,s in Lean (2021
    versions) \cite{kovacs2021smalltt}.
  \item Background: Lean's height-ordered unfolding applies only when neither side has
    metavariables \cite{lean2026defeq}; \code{@[elab\_as\_elim]} infers a motive from the
    expected type \cite{lean2026elabsrc}; pairs of values preserve unfoldings
    \cite{kovacs2024wits}; Rocq ignores mismatched types after postponing and relies on the
    final re-check \cite{kovacs2026wits}; expression-level \code{unfold} complements opaque
    definitions \cite{gratzer2025unfolding}.
\end{itemize}

\subsection{What Stitchwort could take}
\begin{enumerate}
  \item One declaration-level postponement queue with Lean's ladder (T1, D1990): now in
        design as SW-817.
  \item Document and test the lossy rule for stuck calls (T2): small, no behaviour change.
  \item SW-805 justified as a name-resolution rule, with an ambiguity message naming both
        types (T3).
  \item smalltt's three benchmarks as elaboration probes (T4).
  \item Motive inference for \code{T.induct} applications (T5).
  \item Messages say which expression solved an unknown; never a bare \code{?n} (T6).
  \item Keep both named and forced expected types in \code{check\_inner} (SW-858).
  \item Expression-level \code{unfold f in e} (SW-796, SW-798).
  \item Lean's scope-argument constant approximation (measure first, SW-815).
  \item Pattern-matching unification \cite{richter2026pmu}: research, not now.
\end{enumerate}
Not recommended: functions-as-constructors unification \cite{libal2021fcu}, complete search
\cite{norman2026canonical}, contextual typing \cite{xue2024contextual}, an
extensible-constraints architecture, and tabled resolution \cite{selsam2020tabled}, which only
matters with overlapping instances.

\subsection{Acted on, and open}
D1990 led to \file{docs/proposals/postponement.md}. Its first draft measured that the
queueable \code{E-CANNOT-INFER} sites fire 1,073 times in accepted code, always caught, and
proposed a second pass; D2012 made the end state one queue in the single pass. The draft 2
prototype changed one acceptance and 26 hashes (statements identical), and no error message.
The external review is being worked into draft 3 (D2015). SW-805 is merged as a last resort
(D1991). Open: the note's D2013 and D2014 (Section~\ref{sec:decisions}).

% =====================================================================
\section{Proof automation}\label{sec:pa}

\textbf{Register:} \file{docs/research/proof-automation.md}. 72 screened, 29 carded, 25 read, 4 NOT
READ; redone.

\subsection{Key findings}
\begin{itemize}
  \item \textbf{Equalities among hypotheses.} Three probes fail with \code{cases} today \run:
    two hypotheses equating a list with \code{Cons} and with \code{Nil}; \code{f a = c} from
    \code{a = b} and \code{f b = c}; and a local \code{(y: Nat) -> lookup = Some y -> y < 5}
    applied to \code{lookup = Some y0}. Congruence closure with a proof forest settles the first
    two \cite{selsam2016cc,lean2026grindsrc}; instantiating local $\forall$ facts settles the
    third. Selsam and de Moura's subsingleton rule handles proof arguments, and Stitchwort
    already has the lemmas it needs (axiom K, \code{Bool.uniqueEq}, \code{DecEq.uniqueEq}).
  \item \textbf{Lean's \code{grind}.} All proofs are by contradiction on a shared whiteboard;
    it produces ordinary proof terms; constraint propagation closes goals with no case split;
    case analysis is bounded and ``not exhaustive''; it is ``not designed for goals whose search
    space explodes combinatorially'' \cite{lean2026grind,demoura2025grind}. Statically
    predicting the right split is no better than random: on rescuable split failures stock
    \code{grind} gets 0\%, learned models 30--53\%, uniform random 57\% and one-step lookahead
    100\% \cite{wang2026learned}.
  \item \textbf{The counterexample search misses some false theorems silently.}
    \code{shortIncreasing (xs: List Nat) (h: increasing xs = True): xs.length < 6} is false, yet
    \code{check --counterexamples} prints ``none in 10000 cases'' \run. The verdict is ``gave
    up'' only when hypotheses hold in under 10\% of inputs
    (\file{crates/driver/src/counterexample.rs:82}); here they hold often, but only for short
    lists. Recent work builds generators from the hypothesis itself: Palamedes synthesises them
    from Lean \code{Bool} predicates (32 benchmarks; 80\% synthesised in under 3\,s)
    \cite{goldstein2026palamedes}; QuickChick orders hypotheses by ``density''
    \cite{mittelman2026testing}; derived checkers are fuel-bounded semi-decision procedures,
    not \code{Decidable} instances \cite{paraskevopoulou2022computing}. Structural shrinking, the
    kind Stitchwort has, is competitive \cite{keles2026shrinking}; Plausible filters as
    Stitchwort does \cite{plausible2026}.
  \item \textbf{Case theorems from the definition.} Lean's \code{f.fun\_cases} and
    \code{f.induct} add an equation for each discriminant, ``just as if the user wrote
    \code{match h : x with \ldots}'', and are generated on first use \cite{lean2026funind}.
  \item \textbf{Show what search found.} Sledgehammer writes a one-line proof with
    instantiated facts in 70.0\% of cases \cite{bartl2025instantiations}; querysmt writes a
    self-contained script so the proof no longer needs the solver \cite{clune2026hint};
    APOLLO's loop tries built-in tactics at each \code{sorry} \cite{ospanov2025apollo}.
  \item \textbf{External solvers are expensive.} lean-smt supports about 200 of cvc5's 662
    proof rules and verifies 15,271 of 21,595 proofs (71\%) \cite{mohamed2025leansmt};
    Lean-auto reconstructs only through Duper \cite{qian2025leanauto,clune2024duper}; crush
    closes 693 of 754 conditions (91.9\%) with reconstruction against \code{grind}'s 641
    (85.0\%) \cite{he2026crush}; premise selection adds 21\% more goals
    \cite{zhu2026premise}.
  \item Background: small proofs from congruence closure \cite{flatt2022small}; Lean's
    \code{omega} \cite{lean2026omega} and termination guessing \cite{lean2026guesslex};
    reflective rewriting scales where proof-producing rewriting is quadratic in binders
    \cite{gross2022rewriting}; e-graph rewriting under binders is hard enough that lean-egg's
    authors moved to \code{grind} \cite{rossel2024egraphs}; keyed matching for reflective
    \code{ring} \cite{sakaguchi2022reflexive}.
\end{itemize}

\subsection{What Stitchwort could take}
\begin{enumerate}
  \item A congruence-closure closer in \code{cases} (T1, D2000).
  \item Report ``hypotheses held only up to size $k$'' in the counterexample search, and measure
        it over std and aoc (T2; SW-903).
  \item Local $\forall$ hypotheses instantiated once against the classes (stage 2 of 1).
  \item Generated \code{f.cases} (then \code{f.induct}) on demand (T4, D2001).
  \item Lookahead over other atoms before \code{E-CASES-BUDGET} (T6).
  \item Opt-in reports of what \code{cases} and \code{omega} used; ``try automation at each
        \code{?}'' (T5).
  \item Hypothesis-directed inputs, only if 2 finds cases (T3).
  \item Propagation over \code{Bool} connectives before splitting.
  \item Suggest a \code{decreasing} measure on a termination failure.
  \item Reflective rewriting, if \code{simp [rules]} is ever built.
\end{enumerate}
Not recommended: SMT or ATP back ends, e-graph rewriting under binders, and bit-blasting (its
card is NOT READ, and no program needs bit vectors).

\subsection{Acted on, and open}
SW-903 records the silent miss; SW-904 records congruence closure, with Matt's D2000: evidence
from existing proofs first. On \file{decide.md}: \code{grind}'s own splitting is bounded too,
which supports D1971; the literature adds propagation before splitting and lookahead as options.
Open: D2001.

% =====================================================================
\section{Kernel and type theory}\label{sec:kernel}

\textbf{Register:} \file{docs/research/kernel.md}. 62 screened, 27 carded, 26 read, 1 NOT READ;
redone.

\subsection{Key findings}
\begin{itemize}
  \item \textbf{Recursive propositions and proof irrelevance.} If a proposition can be
    eliminated into data and holds a proof of itself in a field, definitional irrelevance lets
    any proof be unfolded once more, and conversion becomes undecidable
    \cite{gilbert2019sprop,felicissimo2026dpi}. Lean is removing such large eliminations
    \cite{lean13851,lean14977}. Stitchwort's \code{E-PROP-ELIM}
    (\file{crates/core/src/check.rs:915-929}) does not ask whether a field mentions the type
    itself, so \file{proof-irrelevance.md}'s \code{Loop} is accepted. Rocq implements
    irrelevance in an untyped conversion with relevance marks, the same idea as SW-93's D1$'$
    masks \cite{gilbert2019sprop}. A cast that reduces on convertible endpoints breaks
    normalisation with an impredicative \code{Prop} \cite{abel2020failure}; Stitchwort's
    \code{J} computes only on \code{rfl}, so this is outside the kernel.
  \item \textbf{Record eta needs a field that isn't a proof.} Once proofs are irrelevant, a
    record of proofs is unit-like and eta on it breaks transitivity; Rocq requires a relevant
    field \cite{lennonbertrand2025certify,rosain2026sort}, and Lean's kernel hit the same case
    \cite{leankernelarena}.
  \item \textbf{R3/R4 have a published confluence argument that needs no termination.} The
    Taming of the Rew uses parallel plus, Stitchwort's four rules, as its example: with a fifth
    ``parallel'' rule they satisfy the triangle property, which gives confluence without
    normalisation \cite{cockx2021taming}. The Rewster shows the naive typing check unsound under
    cumulativity, which Stitchwort lacks \cite{leray2026rewster}.
  \item \textbf{What broke other kernels in 2026.} Rocq's guard checker (CVE-2026-72703, 72705)
    \cite{rocqcve2026}; Lean's order-dependent cache \cite{lean14806} and name-trusted native
    \code{Nat.add} \cite{lean15908}; the Arena's forged \code{Quot.ind}
    \cite{leankernelarena}. An agent team found seven proofs of \code{False} and three other bugs
    in Rocq, none in the official Lean kernel (four other bugs), two in nanoda and one in
    lean4lean \cite{sterin2026falsehood}. Stitchwort's natives are bound by name and checked by
    type and arity only (\file{crates/core/src/native.rs:379-455}), the configuration Lean4Lean
    calls ``technically unsound'' \cite{carneiro2026lean4lean}. Lean4Lean also reports that
    soundness bugs occur where ``the developer did not even consider the error condition''.
  \item \textbf{Quotients and erasure.} Danielsson proves that with quotients and erased
    transport, closed programs' extractions compute the right numeral even where source
    reduction is stuck (Theorem 11.3) \cite{danielsson2026erased}: the nearest published result
    to \file{map-equality.md} \S7.4.5's ``intended invariant'', under different assumptions.
  \item \textbf{Mechanised checkers and their cost.} Lean4Lean runs 1.21--1.40 times slower
    than Lean's kernel \cite{carneiro2026lean4lean}; MetaCoq's development is over 300\,kLoC and
    found one kernel incompleteness, not a soundness bug \cite{sozeau2025metacoq};
    Martin-L\"of \`a la Coq is about 30\,kLoC \cite{adjedj2024mlacoq}.
  \item Not recommended now, each sound research no ticket needs: observational equality
    \cite{pujet2022obs,pujet2025cicobs}, universe polymorphism and cumulativity
    \cite{geran2026levels,bezem2023universe,sterbac2026fussfree}, a concurrent convertibility
    checker \cite{courant2026lazy}, definitional inversion without normalisation
    \cite{carneiro2026definv}, typed confluence \cite{felicissimo2026confluence}.
\end{itemize}

\subsection{What Stitchwort could take}
\begin{enumerate}
  \item No large elimination from recursive propositions, measured first (T2, D1980).
  \item Must-reject tests for the 2026 shapes (T1) \emph{(done, 14 programs)}.
  \item Record eta needs a non-proof field (T3, D1981).
  \item \S9.4 on the triangle property (T4) \emph{(done)}.
  \item Natives: check the defining equations at binding, or pin by hash (T5; SW-898).
  \item One sentence in \file{docs/core.md}: memo tables are keyed by exact pairs, never closed
        transitively (T6). Stitchwort's memos already are: \code{CallMemo}
        (\file{crates/core/src/value.rs:2334}) and the equal-pairs record of
        \code{same\_stuck\_trees} (\file{value.rs:2736}).
  \item Cite Danielsson's Theorem 11.3 in \file{map-equality.md} (T6).
  \item Courant and Leroy's conversion probes (T7).
  \item Inputs to a future universe-polymorphism note (research, not now).
  \item Observational equality for funext and recursion through quotients (research, not now).
\end{enumerate}

\subsection{Acted on, and open}
T1 and T4 are done; T1 found SW-902. D1980 and D1981 are decided and recorded on the tickets
they constrain. SW-898 (pin the natives) is in the backlog. No kernel decision is open.

% =====================================================================
\section{Inductive types and recursion}\label{sec:inductives}

\textbf{Register:} \file{docs/research/inductives.md}. 56 screened, 21 carded, all read; audited.

\subsection{Key findings}
\begin{itemize}
  \item \textbf{Stitchwort's big choices are confirmed.} Lean removed reducible well-founded
    recursion in 4.25 (\#10714) \cite{leanreleases}, and Leroy's pearl supports recursion on a
    measure with proofs by induction on the argument \cite{leroy2024wf}. Sized types remain
    impractical: backward-compatible size inference slowed one file of Coq's standard library,
    \code{MSetList}, from 15.1\,s to 83.7\,s (5.5$\times$), ``intrinsic and unavoidable''
    \cite{chan2023sized}; Agda's sized types are inconsistent because Agda proves
    $\infty < \infty$ \cite{laarakker2026large}.
  \item \textbf{Generated induction and case principles.} Lean's \code{f.induct} and
    \code{f.fun\_cases} follow a function's branches, with one hypothesis per recursive call
    \cite{lean2026funind}. Breitner's warning: facts that exist only inside a termination proof
    aren't available to the induction principle (``The cake is a lie'')
    \cite{breitner2025extrinsic}. Generated definitions checked afterwards are Stitchwort's
    route already; Ko et al.\ show generation costs far more than the definitions
    themselves (about 66\,s against 5\,s) \cite{ko2022reflection,escot2022generic}.
  \item \textbf{Automatic \code{attach}} rewrites \code{xs.map f} to its attached form so a
    recursive call inside \code{map} gets its membership fact \cite{lean2026preprocess}.
  \item \textbf{Termination-measure suggestions.} GuessLex tries candidate measures and
    prints the one that works, or a table per call \cite{lean2026guesslex}.
  \item \textbf{Codata.} Checking termination and productivity separately ``is not enough''
    once inductive and coinductive types nest \cite{hyvernat2025sizechange}; Lean added only
    coinductive predicates \cite{leanrecursive}.
  \item Smaller points: Lean's \code{partial\_fixpoint} reasons about possibly non-terminating
    functions \cite{leanrecursive}; derived enumerations in linear size
    \cite{mathlibderivefintype}; derived checkers for relations \cite{paraskevopoulou2022computing};
    representation of \code{Nat}-like families \cite{theocharis2025custom,idris2builtins};
    constructors resolved by expected type \cite{leandottedident,agda280changelog}; motive
    inference by pattern-matching unification \cite{richter2026pmu}.
\end{itemize}

\subsection{What Stitchwort could take}
\begin{enumerate}
  \item Generated \code{f.cases} and \code{f.induct} on \file{derive_induct.rs}'s pattern (T1,
        D1910); \code{T.induct} itself is done.
  \item \sout{Derived \code{T.all} for enumerations}: superseded by \file{decide.md}'s D1972.
  \item \code{E-TERMINATION} suggests a \code{decreasing} clause (T3).
  \item Automatic \code{attach} for calls inside list callbacks (T2).
  \item Constructor short form resolved by expected type (T6, SW-510).
  \item Audit runtime representations of \code{Fin}, \code{Vec}, \code{Member}, \code{Mod} (T5).
  \item Record that sized types stay closed (SW-199 option D).
  \item Motive inference by pattern-matching unification as a fallback (research).
  \item Derived decision procedures and generators for inductive relations (research).
  \item Partial correctness for \code{Partial} functions; coinductive propositions (research,
        not now).
\end{enumerate}

\subsection{Acted on, and open}
The audit withdrew D1912: it rested on day~21 needing enumerations only, but its
\code{DirectionKey} has a field. \code{T.induct} (D1823) and \code{auto}'s proof of
\code{decreasing T.size p} from an attached element (D1450's ask) were built in nested types
stage~3. Open: D1910 and D1911.

% =====================================================================
\section{Infrastructure and trust}\label{sec:infra}

\textbf{Register:} \file{docs/research/infrastructure.md}. 68 screened, 29 carded, 28 read, 1 NOT
READ (Eris); audited.

\subsection{Key findings}
\begin{itemize}
  \item \textbf{Oracles from the corpus.} Hephaestus found 156 typing bugs (137 confirmed, 85
    fixed) in javac, kotlinc and groovyc; 74 of them (52 inference, 22 soundness) came from two
    mutations of programs that already compile, which its generator alone missed: delete type
    annotations (must still compile) and replace a type by an incompatible one (must not)
    \cite{chaliasos2022hephaestus}. Thalia builds programs from library signatures and found 84
    bugs, at least 42 of them missed by Hephaestus \cite{sotiropoulos2024thalia}. Stitchwort's
    elaborator fuzzer has neither oracle. Content hashes make the erasure oracle stronger here:
    removing an inferable annotation must leave the hash unchanged.
  \item \textbf{Generators.} Generate at the IR the backends share, as Rustlantis does with
    MIR (22 bugs) \cite{wang2024rustlantis}; choose types non-locally so functions use their
    arguments (four times faster to find GHC bugs) \cite{frank2024useless}; score generators by
    mutants \cite{shi2023etna}; a second implementation is an oracle
    \cite{sharma2023rustsmith}; an LLM can seed inputs for an unfamiliar language
    \cite{xia2024fuzz4all}.
  \item \textbf{Keep the reference evaluator in every comparison.} Dafny's compiler fuzzers
    found 24 new bugs, 9 of them soundness bugs; bugs in a shared front end make ``all
    back-ends \ldots\ exhibit the same incorrect result'', so they added an independent
    interpreter \cite{donaldson2024fuzzd}. Stitchwort's runtime and JS backend share their
    erasure decisions (\file{crates/rt/src/erasure.rs}), while \code{run --reference} uses the
    kernel's own erasure.
  \item \textbf{Fuzz the trusted base, not the proofs.} 105,823,818 executions against a
    zlib proved in Lean found no bug in the verified code, one heap overflow in Lean's C++
    runtime and one crash in an unproved parser \cite{kirancodes2026leanzip}.
  \item \textbf{Cache keys.} Lean keeps proofs private to their module, so proof edits don't
    rebuild importers \cite{leanmodules}; build-system vocabulary \cite{mokhov2020bsalc};
    durability levels \cite{rustanalyzer2023durable}; coarse roots invalidate everything
    \cite{mathlibcache2023}; incremental elaboration within declarations \cite{leanreleases}.
    Stitchwort's tests run erased (\file{crates/driver/src/driver.rs:1241-1243}), so a test
    result keyed on the erased program can't be changed by a proof edit.
  \item \textbf{Agent-driven development.} ``The task verifier is nearly perfect, otherwise
    Claude will solve the wrong problem'' \cite{carlini2026compiler}. Agents under pressure edit
    tests; a ``flag for human intervention'' exit cut GPT-5's cheating from 54\% to 9\%
    \cite{zhong2025impossiblebench}. Detailed type errors help agents fix code
    \cite{krishnamurthi2026ablation}, and compiler feedback took GPT-5 from 22 to 54 of 56 Idris
    exercises \cite{li2026idris}. Mathlib's bot lists renamed lemmas and debt metrics on each
    pull request \cite{baanen2025mathlib}; Unison's renames are ``same hash, new name''
    \cite{unison2026bigidea}.
  \item \textbf{Editors and machines.} Type definitions are the most useful context for a
    model filling a hole \cite{blinn2024chatlsp}; LSP is the wrong interface for machines
    \cite{aniva2025pantograph}; presentations that keep a reference to each subterm
    \cite{nawrocki2023ui}; check what is on screen first \cite{coqlsp}; total error recovery
    \cite{zhao2024marking}.
  \item Also: MutDafny mutates programs to find weak specifications (on average one in every 241 lines)
    \cite{amaral2026mutdafny}; Sisyphus repairs proofs, cheap only in its first, differential
    step \cite{gopinathan2023sisyphus}; verified extraction is safe only for first-order
    interop, which supports Stitchwort's first-order JS exports \cite{forster2024extraction}.
\end{itemize}

\subsection{What Stitchwort could take}
\begin{enumerate}
  \item A lane trust diff: changed statements, renames, must-reject, snapshot and budget
        changes, axioms (T1, D1930).
  \item Erase and overwrite oracles in \file{fuzz_elab.rs} (T2; SW-686, SW-687).
  \item A kernel type-preservation sweep: evaluate, read back, re-check (T3).
  \item Rotate the sampled sweeps' offset (T4).
  \item Score fuzzers against missed \code{cargo mutants} mutants (T5, D1932).
  \item Weak-theorem detection by mutating functions (T6).
  \item A detailed error mode for agent lanes (T7).
  \item Goals list the data types they mention (T8).
  \item SW-688 as an IR-level backend generator, always compared with \code{run --reference}
        (T9).
  \item Test results keyed on the erased program (T10, D2030).
  \item Measure cross-namespace proof unfolding before deciding opaque theorems (T11, D1931).
  \item Deprecated aliases for std renames (T12).
\end{enumerate}

\subsection{Acted on, and open}
The audit corrected Hephaestus's numbers (the 29 October card had the anonymous submission
copy's) and found that tests run erased, which simplified T10 and produced D2030. Open: D1930,
D1931, D1932, D2030. Eris is for Matt to fetch (Section~\ref{sec:limits}).

% =====================================================================
\section{Compilation and runtime}\label{sec:comp}

\textbf{Register:} \file{docs/research/compilation.md}. 52 screened, 23 carded, 22 read, 1 NOT READ;
audited.

\subsection{Key findings}
\begin{itemize}
  \item \textbf{Most standard ideas are already in.} The October runtime survey's ideas
    (erased IR, register bytecode, reuse tokens, small numbers inline, TRMC, fused state
    handlers, Cranelift) are built.
  \item \textbf{Arrays with in-place update.} Lean's \code{Array.write} updates destructively
    when unshared \cite{ullrich2019beans}; a record of an efficient encoding plus an erased
    invariant and a view is the general method \cite{allais2023builtin}; FP$^2$'s dynamic
    uniqueness test costs about 1\% when nothing is reused \cite{lorenzen2023fp2}. The risk:
    in Huisinga's own \code{Array.groupBy} example, the array also sits in a map, so every
    \code{push} copies it (``accidentally quadratic''); the thesis presents this as its own
    illustration, not a bug in Lean's library \cite{huisinga2023uniqueness}.
  \item \textbf{Borrowing has costs.} Inferring borrows removes 75--100\% of increments on 11 of
    13 benchmarks against Perceus without reuse, 1.48$\times$ geomean, but blocked 6.4\% of
    in-place updates in one benchmark \cite{brandon2026borrows}. Koka declined Lean's borrow
    inference as ``not safe for space'' \cite{lorenzen2022frame}. In Lean's own ablation,
    removing borrowing made \code{const\_fold} faster (0.90) \cite{ullrich2019beans}; Lean now
    lets users annotate borrows, with tail calls taking priority \cite{leanreleases,leaninferborrow}.
  \item \textbf{Forced indices.} A constructor argument that appears in its result indices
    needn't be stored, saving 10--20\% of memory on vectors \cite{brady2004indices}; Agda's GHC
    backend drops erased fields \cite{agdaruntimeirrelevance}. Stitchwort's \code{VCons} stores
    its length, and \file{vec.md} measured 5--10\% per operation.
  \item \textbf{Dispatch is not the lever.} Tail-call dispatch in Rust ran 1.5--2$\times$ faster
    natively but 1.2--4.6$\times$ slower in WebAssembly engines, and needs nightly
    \cite{keeter2026tailcall}; CPython's version gained 1--5\% against a good baseline
    \cite{elhage2025tailcall,jin2025sorry}; Wizard's threaded dispatch, 14\% on average
    \cite{titzer2022inplace}; Deegen has no ablation \cite{xu2026deegen}; register bytecode was
    wasmi's big step \cite{wasmi032}. Dispatch is about 19\% of Stitchwort's machine time
    (\file{native-compilation.md} \S1.2), so SW-813 should start from the call path.
  \item \textbf{Small numbers in JS.} The JS backend makes every \code{Nat} a \code{bigint};
    day~2's brute force took 66.8\,s under \code{--js} against 23.1\,s on the machine. Scala.js
    measured \code{BigInt} longs 60 times slower than its own, in 2018 \cite{tc39bigint117}.
  \item \textbf{Measure with instruction counts} on a loaded machine \cite{instrcounts}.
  \item \textbf{Corner cases in shifts.} Verifying CertiCoq-Wasm found that WasmCert modelled Wasm
    shifts like C shifts, and a corner case that led to an unsoundness in Rocq's
    \code{vm\_compute}
    \cite{meier2025certicoqwasm}.
  \item Not transferable now: lazy reference counting, which bounds pauses rather than saving
    time \cite{lam2024lazy}; a Wasm backend, which still needs its own stack for deep recursion
    \cite{wasmweb2024}. Context: lambda set specialisation is the general form of
    \file{opt.rs}'s known-argument copying \cite{brandon2023lss}; zipper-based tree updates
    \cite{lorenzen2024fiptree}.
\end{itemize}

\subsection{What Stitchwort could take}
\begin{enumerate}
  \item Benchmarks record retired instructions and bytecode counts beside time, and gate on
        instructions (T1; SW-659).
  \item Shift and bitwise natives tested at and past the word width in every runner (T7;
        SW-808).
  \item Count reference-count operations per run, the gate for any borrowing (T2, D1921).
  \item \code{Array} specified as a list, in place when unique, with a copy counter (T6, D1920;
        SW-112).
  \item SW-813: start from the call path; no tail-call dispatch while the playground is Wasm
        (T4).
  \item Check reuse-token placement against Koka's fragility examples (T8).
  \item JS \code{Nat} as \code{number} below $2^{53}$, \code{bigint} above, measured first (T5,
        D1922).
  \item A ``runs in place'' report for std functions.
  \item \code{Map} insert as a zipper loop that reuses every node (SW-814).
  \item Borrowed parameters, only if item 3 shows a large share.
  \item Don't store forced indices; drop self-passed relevant arguments.
  \item A size counter for lowered code.
  \item Unboxed scalar fields in cells, scalars after pointers.
\end{enumerate}

\subsection{Acted on, and open}
The audit corrected C8 and two file references; nothing has been built yet. Open: D1920, D1921,
D1922.

% =====================================================================
\section{Cross-cutting themes}\label{sec:themes}

Each theme below appears in at least two registers; none is new to this summary.

\subsection{Soundness hardening, and where real bugs were found}
The kernel register's 2026 table (\S5.4) was the most productive single exercise. Three of its
shapes are closed by construction in Stitchwort and two by design, which the 14 tests confirmed;
one (natives bound by name) is a known small gap. The bit-level shortcut had not been audited,
and porting it, Lean4Lean's overflow shape, found SW-902. That matches Lean4Lean's observation
that bugs live where nobody considered the error condition \cite{carneiro2026lean4lean}, and
Stérin's that adversarial agents given past bugs find new ones \cite{sterin2026falsehood}. The
natives gap is SW-898. Two hazards were found before the
features that would expose them: recursive propositions under proof irrelevance (D1980) and
record eta over proof-only records (D1981). The infrastructure register adds the testing side:
must-reject examples paired with accepted neighbours, a type-preservation sweep, and an
``either'' category for honest corner cases \cite{leankernelarena}.

\subsection{Proof automation a small, predictable kernel can afford}
Every automation register finds the same shape: search outside the kernel, produce an ordinary
proof term or certificate, and let the kernel check it. \code{grind} does this and reads only
what it has collected \cite{lean2026grind}; Stitchwort's \code{omega}, \code{ring} and
\code{cases} already do. External solvers buy a few points at the cost of large
reconstruction libraries \cite{mohamed2025leansmt,he2026crush}, and Stitchwort's failures are
about facts in scope, not large lemma libraries. Proposals that keep the kernel unchanged
(congruence closure in \code{cases}, generated case theorems, write-back of found proofs) come
first; the one proposal that would trust more, monad laws as conversion rules, is a design note
only (D1850).

\subsection{Elaboration order and postponement}
Three registers meet here. Elaboration recommends one postponement queue
\cite{lean2026elabsrc,obrien2025omniml}; inductives and proof automation recommend motive
inference from the expected type (\code{elab\_as\_elim}, pattern-matching unification
\cite{richter2026pmu}); SW-805's field-name rule is acceptable because it depends only on the
text, not on inference order \cite{gundry2020ghc366}. The common principle: anything that depends
on inference order must wait in one place, with a fixed ladder of defaults, and anything decided
from the text alone can run at once.

\subsection{Testing oracles}
Differential testing against an independent implementation is the oracle the literature trusts
most: Rustlantis compares backends \cite{wang2024rustlantis}, fuzz-d adds an interpreter
because shared code fools back-end comparisons \cite{donaldson2024fuzzd}, and Sisyphus runs old
and new versions on the same inputs \cite{gopinathan2023sisyphus}. Stitchwort's reference
evaluator is that implementation and must stay in every comparison. Mutation gives oracles
without a reference: Hephaestus's erase and overwrite \cite{chaliasos2022hephaestus}, MutDafny's
weak-specification check \cite{amaral2026mutdafny}, ETNA's scoring of generators
\cite{shi2023etna}. The counterexample search is itself an oracle with a blind spot (SW-903), and
the remedy is the same as for any oracle: report what it did not reach.

\subsection{Trusted base}
The registers keep returning to what is trusted: the kernel's re-check is the soundness
argument for the lossy unifier; natives are trusted and should be pinned to the definitions
they replace; R3/R4 are trusted conversion rules with a published confluence argument; the
runtime, JS codecs and natives are where fuzzing finds bugs \cite{kirancodes2026leanzip}; JS
exports stay first-order because callbacks can be impure \cite{forster2024extraction}; an erased
test key would add trusted cache code, guarded by a differential test (D2030).

\subsection{Measure before designing}
Many recommendations are measurements: count reference-count operations before borrowing
(D1921), time JS \code{number} against \code{bigint} (D1922), measure how often hypotheses hold
only for small inputs (SW-903), survey existing proofs before congruence closure (D2000), count
cross-namespace proof unfoldings before opaque theorems (D1931), and the postponement note's
counter of queueable sites (1,073) before its design.

% =====================================================================
\section{Recommendations}\label{sec:recs}

Table~\ref{tab:recs} merges the seven registers' rankings into one list. Priority weighs value
(soundness first, then open tickets and friction in real programs) against effort, as the
registers do; within a priority band the order follows the registers. ``Open'' means proposed
and not started; done items are kept for the record.

\begin{longtable}{@{}rL{5.2cm}L{2.45cm}L{1.85cm}L{2.4cm}L{1.6cm}@{}}
\caption{One prioritised list across all clusters.}\label{tab:recs}\\
\toprule
\# & Item & Value & Effort & Ticket / decision & Status \\
\midrule
\endfirsthead
\toprule
\# & Item & Value & Effort & Ticket / decision & Status \\
\midrule
\endhead
\bottomrule
\endfoot
1 & Must-reject tests for 2026 kernel-bug shapes (kernel T1) & soundness & small & D1983 & done \\
2 & Fix universe-level overflow & soundness & small & SW-902, D1984 & done \\
3 & \S9.4 on the triangle property (kernel T4) & soundness argument & doc & D1983 & done \\
4 & Refuse large elimination of recursive propositions (kernel T2) & soundness before SW-93 & small kernel & D1980; SW-93 & decided, with SW-93 \\
5 & Record eta only with a non-proof field (kernel T3) & soundness after SW-93 & doc & D1981; SW-55, SW-789 & decided (reminders) \\
6 & Pin natives by hash or check their equations (kernel T5) & trusted base & small kernel & SW-898 & open \\
7 & Lane trust diff (infra T1) & review of agent work & small tooling & D1930 & open \\
8 & Counterexample search reports the size reached (PA T2) & silent misses & small & SW-903 & open \\
9 & Erase and overwrite oracles (infra T2) & elaborator soundness & tooling & SW-686, SW-687; D1932 & open \\
10 & Kernel type-preservation sweep (infra T3) & kernel soundness & small & --- & open \\
11 & One postponement queue (elab T1) & SW-817, 847, 885, 509 & elaborator & D1990, D2012, D2013, D2014; SW-817 & in progress \\
12 & SW-805 field-name rule as name resolution (elab T3) & SW-805, SW-817 & small & D1991; SW-805 & done (rule) \\
13 & Benchmarks gated on instruction counts (comp T1) & every speed lane & small & SW-659 & open \\
14 & Shift natives at and past the word width (comp T7) & correctness & small tests & SW-808 & open (with SW-808) \\
15 & Document and test the lossy rule (elab T2) & predictability & small & --- & open \\
16 & Score fuzzers with missed mutants (infra T5) & choosing SW-688/689 work & small & D1932 & open \\
17 & Congruence-closure closer in \code{cases} (PA T1) & SW-608 & elaborator & D2000; SW-904 & evidence first \\
18 & Generated \code{f.cases} and \code{f.induct} (PA T4, inductives T1) & SW-802, SW-809, SW-199 & elaborator & D2001, D1910 & open \\
19 & Count reference-count operations (comp T2) & gate for borrowing & small & D1921 & open \\
20 & Test cache keyed on the erased program (infra T10) & warm tests after proof edits & small & D2030 & open \\
21 & Rotating sample offset (infra T4) & coverage & small & --- & open \\
22 & Handler fast-path report (effects T5), after re-measuring SW-861 & SW-861, SW-791 & small & --- & open \\
23 & \code{E-TERMINATION} suggests a \code{decreasing} clause (inductives T3) & messages & small & --- & open \\
24 & Lookahead before \code{E-CASES-BUDGET} (PA T6) & SW-608 & small & --- & open \\
25 & Show what automation found; try automation at each \code{?} (PA T5) & readability, agents & small & --- & open \\
26 & smalltt probes; Courant--Leroy probes (elab T4, kernel T7) & proof-check speed & small & --- & open \\
27 & \code{decide}: extend \code{cases}, derive \code{Finite} & SW-800, SW-807 & elaborator & D1971, D1972, D1973 & awaiting review \\
28 & Handler specifications as predicate transformers (effects T4) & SW-570 & design note & D1851 & open \\
29 & \code{Array} specified as a list, in place when unique (comp T6) & SW-112 & design, runtime & D1920 & open \\
30 & JS \code{Nat} as \code{number}, measured first (comp T5) & \code{run --js} & measure & D1922 & open \\
31 & Automatic \code{attach} (inductives T2) & SW-199 & elaborator & --- & open \\
32 & Agent detail mode; goals list their data types (infra T7, T8) & agents & small & --- & open \\
33 & Monad laws as trusted conversion rules (effects T3) & SW-548 & kernel & D1850 & design note only \\
34 & Cross-namespace proof unfolding, then opaque theorems (infra T11) & early cutoff & measure & D1931 & not yet \\
35 & Codata, sized types, observational equality, universe polymorphism & --- & research & D1911 & not recommended \\
36 & SMT or ATP back ends; e-graph rewriting under binders & --- & research & --- & not recommended \\
37 & Tail-call dispatch while the playground is Wasm; a Wasm backend & --- & runtime & --- & not recommended \\
\end{longtable}

% =====================================================================
\section{Open decisions for Matt}\label{sec:decisions}

Table~\ref{tab:decisions} lists every decision still open in the registers and design notes,
with each register's recommendation and its assumptions in full. Each was checked against the
board and \code{git log} on 30 October; resolved ones are left out (D1852, D1912, D1980, D1981,
D1983, D1984, D1990, D1991, D2000 and D2012; the postponement note's D2010 was superseded by
D2012, and its D2011 is kept in draft 2 and not listed there as open). D2013 and D2014 come from
\file{docs/proposals/postponement.md}, which draft 3 (D2015) is revising after its external
review, so their wording may change. D1971--D1973 come from \file{docs/proposals/decide.md},
whose SW-800 awaits a review artefact.

{\small
\begin{longtable}{@{}L{1.1cm}L{4.1cm}L{2.5cm}L{7.15cm}@{}}
\caption{Open decisions.}\label{tab:decisions}\\
\toprule
ID & Question & Recommended & Assumptions, in full \\
\midrule
\endfirsthead
\toprule
ID & Question & Recommended & Assumptions, in full \\
\midrule
\endhead
\bottomrule
\endfoot
D1850 & Should the monad laws of \code{Effect.bind} and \code{Effect.foldComp} hold by conversion
(trusted kernel rules, as \code{Nat.add}'s two extra rules), though their equations need funext? &
Yes to a design note (T3) only; decide the kernel change after it. &
The decision line states none. The deep read (\S4.1) names one: rules whose equations need
funext are acceptable as trusted conversion (some theorems would stop listing \code{Quot.sound}
in \code{stitchwort axioms}; consistency should be safe, because funext holds in the models
Stitchwort's theory already assumes \inferred). Its risks: a normalisation argument for the
whole system must be written out; the elaborator's unification must apply the same rules. \\
\addlinespace
D1851 & A specification layer for handlers in std (\code{wp} per handler, after Lean's
\code{Std.Do}), aimed at SW-570 and rogue's loop? &
Yes, design note first (T4), after T2 lands. &
The decision line states none. The deep read (\S4.3) names one: the cost is in symbolic
evaluation of handlers, not in the property (rogue's notes support it). Its risks: a \code{wp}
that recurses through a thunk will need a data form, as \code{AlwaysTree} did; partial programs
still need fuel. \\
\addlinespace
D1910 & Write one design note covering \code{f.cases} and \code{f.induct}, replacing D231--D233's
``not now'' and answering D1439? &
Yes. &
(A1) the elaborator keeps, or can recompute, the pattern variables, the index equations learned
and the \code{if} conditions at each leaf of a compiled case tree; (A3) generated theorems are
made only when a program names them. \\
\addlinespace
D1911 & Keep codata out of Stitchwort for now (no streams, no coinductive types or propositions),
and revisit only when a program needs infinite data? &
Yes. &
(P1) no planned program needs it; (P2) the \code{Partial} model stays a one-element set.
Hyvernat's result means a later design needs its own totality checker, not a second size-change
pass. \\
\addlinespace
D1920 & For SW-112, write the design note around a list-specified \code{Array} with in-place
update on unique values in the machine and native code, rather than a persistent tree or a
uniqueness type? &
Yes: design note T6; build only after it settles the JS choice. &
Users thread arrays linearly most of the time; a copy counter makes the other case visible; JS
may be slower. \\
\addlinespace
D1921 & Spend a small lane counting reference-count operations (T2), and design borrowed
parameters only if the count shows a large share? &
Yes. &
The count is cheap to add next to \code{alloc\_trace}. \\
\addlinespace
D1922 & Measure small \code{Nat}s as JS \code{number}s (T5) now, or leave \code{--js} as it is
because native code is the CLI's speed route? &
Measure now: JS is the faster runner on parse-heavy programs (SW-593). &
The 2018 \code{BigInt} gap hasn't closed; T5 tests exactly that. \\
\addlinespace
D1930 & Make the trust diff (T1) a standard part of reviewing every lane? &
Yes: build T1, then add ``read the trust diff first'' to \file{docs/workflow.md}. &
Agents keep writing most code and proofs; reviewing, not writing, is the bottleneck; the list is
usually empty. \\
\addlinespace
D1931 & Should theorems be opaque to other namespaces, so a proof edit stops at its namespace
(Lean's module-system rule)? &
Not yet: run T11 and T10 first, decide on the numbers. &
T11 finds few cross-namespace proof unfoldings; std lemma edits are common enough to matter; an
explicit opt-out like Lean's \code{@[expose]} is acceptable. \\
\addlinespace
D1932 & For SW-686/687/688, do the corpus oracles (T2, T3) and the mutant scoreboard (T5) before
building any generator? &
Yes. &
Mutating Stitchwort's large example and AoC corpus reaches more code per hour than a new
generator, as the Rust results suggest. \\
\addlinespace
D2030 & Key test results on the erased program (T10), so editing a proof never re-runs a test? &
Yes, with the differential test; first measure how many tests a typical std proof edit re-runs. &
(1) every test runs erased, on the fast runtime's erased IR or on the evaluator with erasure
switched on (\file{crates/driver/src/driver.rs:1229-1243}; tests of partial code too,
\file{crates/driver/src/run.rs:545}), so no proof body or type argument is computed and a proof
edit can't change a test's outcome; (2) the erased hash of a test can be computed from
\file{erase.rs}'s decisions, with each reference hashed by that declaration's own erased hash,
and it changes whenever anything that runs changes, including a definition's body, a data
type's constructors and the evaluation modes already in the key; (3) a second hash per
declaration in trusted cache code is acceptable if a differential test (the same results with
either key, over the test corpus) guards it, as \file{check-cache.md} \S5.4 guards the check
cache; (4) the saving matters, because std lemma edits re-run many tests today (unmeasured). \\
\addlinespace
D2001 & For SW-802 and SW-809, design generated case theorems (\code{f.cases}, Lean's
\code{fun\_cases}) instead of per-leaf equation lemmas (D231)? &
Yes. &
(F1) the elaborated case tree records which equation each leaf tested, readably enough to state
the branches; (F2) today's \code{cases} does not already prove the day~16 and day~23 statements
that failed with \code{E-MATCH-MOTIVE}; T4 checks this first, and if \code{cases} does, the
decision lapses to ``document \code{cases} for this''. \\
\addlinespace
D2013 & Run SW-805's field-name rule as a default in the queue's ladder, and default number
literals in the same ladder after it, instead of at the literal? &
Yes. &
(R1, from D1991) the field-name rule looks only at the literal's field names and the names in
scope at the literal, never at types learned later; running it later doesn't change what it
looks at, only whether something else decided first. (N1) a number literal that waits ends with
the type the rest of the declaration fixes, and \code{Nat} where nothing does; so in a program
accepted today, where each such literal was made \code{Nat} and everything else agreed with
that, it still ends as \code{Nat}; the change only accepts programs that fail today because a
literal was made \code{Nat} too early; stage 4 tests this over the corpus (inferred, not yet
measured: K4 is not in the prototype). (N2) the field-name rule has priority over the number
default because a record's field types say what its numbers are, never the reverse (\S5.4's
``in a list''). (N3) Lean's numeral default is the precedent: low priority, in the ladder
(\file{Prelude.lean:1411}, \file{SyntheticMVars.lean:237-243}). \\
\addlinespace
D2014 & Resume queued checks at the end of every application, scoped to the items made inside
it, as well as before lifting and at the end of the declaration? &
Yes. &
(S1) Stitchwort's unifier can't wait: a comparison that meets an unfilled placeholder answers
\code{false} at once, and its caller takes another route, so a placeholder must be filled before
the comparisons that would need it; the end of the application that created it is the first
fixed point after the later arguments that may solve its blocker, and before the application's
type is compared with what is expected of it; the prototype shows the end of the declaration is
too late for std (\file{Map/Spec.sw:1862}), and the end of each application is enough for the
whole corpus. (S2) scoping to the
application's own items keeps a failing item from being blamed on, and undone with, an attempt
it doesn't belong to, as Lean's \code{withSynthesize} does; items from outside are resumed at
their own application's end or at the declaration's. (S3) waking inside the unifier would
elaborate syntax in the middle of a comparison that callers inside many of the 44
\code{Snapshot::take} sites branch on; resume points give the same order at fixed places. (S4) if waiting comparisons (\S5.9) are ever added, W1
(the resume point at each application's end) can become less frequent; this decision doesn't
preclude that. \\
\addlinespace
D1971 & Close SW-800 by extending \code{cases} (a decision closer at the leaves, and splitting a
finite variable when nothing is stuck), with no new \code{decide} word? &
Yes. &
\file{decide.md} gives the alternative rather than assumptions: a \code{decide} word evaluating
\code{Decidable} instances, as Lean does, can't do day~21's \code{tabulateSame} and doesn't show
the failing case. The proof-automation register finds that the literature supports it. \\
\addlinespace
D1972 & Build SW-807 as a derived \code{Finite} interface (a list of every value with a
completeness proof) and \code{Finite.check}, for non-recursive constructor types with finite
fields? &
Yes, after D1971's conversion of days 12 and 21. &
Alternative given: enumerations only, as Haskell's \code{Enum}; day~21's \code{DirectionKey} has
a field and would be left out (confirmed by the inductives audit). \\
\addlinespace
D1973 & Never add native evaluation as a proof (\code{native\_decide})? &
Never, until a proof needs more than about 10\,s of kernel evaluation; then revisit with that
proof. &
Alternative given: trust the fast runtime for closed \code{Bool}s, adding it to the trusted base. \\
\end{longtable}
}

D1850 and D1851 state no assumptions in their decision lines; the table gives the assumptions
and risks their deep reads (effects register \S4.1 and \S4.3) name.

% =====================================================================
\section{Limits}\label{sec:limits}

\subsection{Not read}
Seven carded sources were not read. None supports anything in this summary.
\begin{itemize}
  \item Proof automation C26: Morrison and de Moura, \emph{grind: an SMT-inspired tactic for Lean
        4}, I{}JCAR 2026 (no public copy found; its system was read through its manual, source and
        slides).
  \item Proof automation C27: Rossel et al., \emph{Towards Pen-and-Paper-Style Equational
        Reasoning in Interactive Theorem Provers by Equality Saturation}, POPL 2026, doi
        10.1145/3776667 (ACM 403; its tool's README now points to \code{grind}).
  \item Proof automation C28: Bhat, B\"oving et al., \emph{Interactive Bitvector Reasoning using
        Verified Bit-Blasting}, OOPSLA 2025, doi 10.1145/3763167 (no bit vectors yet).
  \item Proof automation C29: Limperg and From, \emph{Aesop}, CPP 2023, doi
        10.1145/3573105.3575671 (no tactic language to search over).
  \item Kernel C27: Liesnikov and Cockx, \emph{Building a Correct-by-Construction Type Checker
        for a Dependently Typed Core Language} (Agda Core), APLAS 2024, doi
        10.1007/978-981-97-8943-6\_4 (TU Delft repository unreachable).
  \item Infrastructure C3: Sotiropoulos and Su, \emph{Enumerating Ill-Typed Programs for Testing
        Type Analyzers} (Eris), PLDI 2026, doi 10.1145/3808320.
  \item Compilation C22: IOG, \emph{Stacks in the JS backend} (2023), a blog post that returns
        404.
\end{itemize}

\subsection{For Matt to fetch}
One paper meets the bar: \textbf{Eris} (PLDI 2026, Distinguished Paper, doi
\href{https://doi.org/10.1145/3808320}{10.1145/3808320}). It is the published form of the
overwrite oracle in recommendation~9, and would shape that design. It is open access at ACM, so
it should open in a browser; every scripted route returned 403 or nothing. Its abstract's figures
(64 bugs, 51 of them unsoundness) are \unverified{} and support nothing here.

\subsection{What wasn't searched}
Each register's \S9 lists what it did not reach. In short:
\begin{itemize}
  \item \textbf{Effects}: Unison's and Koka v3.2's release notes, and the F$^\star$ indexed-effects
        venue, were not checked; the rows they back stay unverified. The 64 screened-only rows
        were not re-screened.
  \item \textbf{Elaboration}: in-window work on elaborating effect rows, incremental elaboration,
        Agda's and Idris~2's elaborator changelogs, RocqPL and Lean workshop programmes.
  \item \textbf{Proof automation}: ITP, CPP, I{}JCAR, CADE and FroCoS programmes year by year;
        Isabelle's Quickcheck and Nitpick after 2022; Agda's and Idris~2's automation.
  \item \textbf{Kernel}: Idris~2's core and its bugs; Agda's kernel bug history; small trusted
        kernels as a topic; hash-consing in production kernels beyond Courant and Leroy.
  \item \textbf{Inductives}: Idris~2's totality changelog, Rocq~9 fixpoints, F$^\star$ and Liquid
        Haskell termination, views papers, non-empty lists, or-patterns; venue sweeps of TyDe,
        CPP, ITP, FSCD and TYPES. One query (Dafny's default \code{decreases}) was never run.
  \item \textbf{Infrastructure}: MetaRocq updates, Agda interface files, F$^\star$ tooling, Nix
        content-addressed derivations, Dhall's pinned imports, Elm's semantic versioning, and the
        CoqPL, Lean workshop and ITP/CPP programmes year by year.
  \item \textbf{Compilation}: the search budget ran out after 36 queries on 29 October, so
        Unison's and Idris~2's backends, MLKit, deforestation, Koka's recent benchmarks,
        persistent vectors, Larose et al.'s AST-versus-bytecode paper, Lorenzen's thesis and Roc's
        docs were not screened.
\end{itemize}

\subsection{What is inferred, not run}
The registers mark their reasoning \inferred. The main inferred claims behind the
recommendations are:
\begin{itemize}
  \item that congruence closure would close probes 1 and 2, and that \code{J}-based congruence
        steps type-check for the dependent calls that matter (D2000's A1--A3);
  \item the critical pairs of the proposed monad-law rules, worked by hand, with one overlap
        (R-fuse against R-assoc) still to check properly;
  \item that \code{f ?a = f x} with \code{f (Succ n) = f n} picks \code{?a := x} (``inferred, not
        run'');
  \item that a record of proofs becomes unit-like under D1$'$ and breaks transitivity;
  \item that Stitchwort's size-change checker would reject \code{f (C2 x) = f (C1 x)}
        (``inferred, not run'');
  \item the expected gains of borrowing (0--10\%, possibly negative), of forced indices and of
        JS \code{number}s, all unmeasured;
  \item that a test key on the erased program changes whenever anything that runs changes
        (D2030 (2)).
\end{itemize}
What \emph{was} run: the proof-automation probes (three \code{cases} failures and the
counterexample miss), the 14 must-reject programs, the parallel-rule test, the universe-overflow
reproduction and fix, and the postponement prototype.

\subsection{This summary}
It cites only works whose cards say ``Full text: read''. Its numbers are taken from the
registers as they stand after the 30 October audits. It did not re-read any paper. Where a
register's file:line reference had moved on \code{main}, the current line is given here (for
example \code{apply\_args\_postponing} is now at \file{check.rs:1903} and \code{CallMemo} at
\file{value.rs:2334}). Bibliographic fields missing from a card were looked up in the arXiv API,
Crossref and DataCite; no DOI was guessed, and entries that couldn't be completed carry a note.

% =====================================================================
\phantomsection
\addcontentsline{toc}{section}{References}
\printbibliography[title={References}]

\end{document}
